\begin{abstract}
Robot world models receive policy commands, while physical transitions are driven by commands transformed by execution software. Hidden controller memory can therefore affect both prediction and the gradients used for policy optimization. We study EAGM, which combines a structured observable-command bottleneck with finite-difference policy-gradient supervision. In a fixed, partially unknown execution-controller setting, independent five-seed experiments support improved physical tracking returns over a matched model without directional supervision, including a simulator-budget control. These benefits do not establish general robot-task transfer. A true end-effector reaching evaluation fails to reproduce the benefit. Frozen-model interventions show that correcting reward geometry alone is insufficient, whereas replacing transition derivatives restores gradient alignment. Subsequent optimizer and start-distribution checks nevertheless fail joint held-out improvement criteria. The results distinguish prediction accuracy, gradient accuracy, and physical update utility, and identify the scope and remaining limitations of execution-aware supervision.
\end{abstract}
\section{Introduction}
A robot command is not necessarily the actuator input that causes the next observed transition. Control software may accumulate target positions, filter incoming commands, buffer actions, or limit command rates. Some of this execution state is observable; other parameters and memories may be unavailable to a deployed predictor. This interface matters when a world model is differentiated to improve a policy, rather than used only to predict a state.

A low prediction error does not constrain all local derivatives. Moreover, even an accurate policy gradient is defined for a particular objective and starting-state distribution. The optimizer then transforms that gradient into an update. Our central question is consequently restricted: can command-structured modeling and simulator-derived gradient supervision improve physical policy returns under partially observed execution dynamics, and where does this approach fail?

EAGM predicts an executed command using a nominal observable update plus a learned residual. A dynamics network maps the physical state and predicted command to the next state. Directional supervision compares a differentiated imagined return with finite differences of a physical rollout at the same policy parameters. The implementation uses a masked, coordinate-normalized Huber objective; it is not a pure cosine loss.

The evidence has three distinct levels. First, paired tracking experiments and independent confirmation seeds support return improvements within the tested controller setting. Second, budget controls and architectural comparisons constrain explanations based solely on extra interaction or an unconstrained auxiliary head. Third, reaching experiments expose a failure boundary and enable frozen-gradient interventions. These interventions are diagnostic privileges, not deployable algorithms.

Our contribution claim is the evaluated combination of an execution-command bottleneck and query-accounted gradient supervision in this restricted setting, together with an explicit analysis of its limitations. We do not claim that gradient-aware model learning is new, that EAGM achieves public-benchmark state of the art, or that an oracle establishes a successful new method. The submitted execution-memory diagnostic is a separate project; the overlap and additional algorithmic evidence must be documented before submission of this manuscript.
